\documentclass{article}
\usepackage{graphicx} % Required for inserting images
\usepackage[a4paper,margin=1in]{geometry}
\usepackage{amsmath, amssymb, float}
\usepackage{enumitem}
\usepackage{fancyhdr}
\usepackage{titlesec}
\usepackage{minted}
\setminted{breaklines}
\usepackage{hyperref}
\usepackage{titlesec}
\usepackage{bm}

\title{knots\_q3\_draft}
\author{Narcis Staicu}
\date{November 2026}

\begin{document}

\maketitle

\section{Introduction}

The \textbf{Kauffman Bracket} is a polynomial invariant \footnote{It is not a knot invariant, as Reidemeister Move I introduces a crossing which alters it} which assings a Laurent polynomial to a knot diagram (not the knot itself). A Laurent polynomial is a polynomial that can contain negative powers (e.g. $A^2 + 3A^{-3}$).

The Kauffman Bracket is written as:

\[
    \langle D \rangle
\]

where D represents the knot diagram.

\hfill

The main idea is to use \textbf{smoothings} to calculate the Kauffman Bracket, instead of calculating it from its complicated knot diagram.

\hfill

A smoothing is a way of removing a crossing from a knot diagram by reconnecting the four strands around the crossing in one of two possible ways, without allowing the strands to cross. 

\hfill

\subsection{Smoothings}

There are two possible smoothings:
\begin{itemize}
    \item A-smoothing
    \item B-smoothing
\end{itemize}

The A-smoothing contribution to the Kauffman Bracket is:

\[
    A \langle A\text{-smoothed diagram} \rangle
\]

The B-smoothing contribution to the Kauffman Bracket is:

\[
    A^{-1} \langle B\text{-smoothed diagram} \rangle
\]

\subsection{Rules}

There are three Kauffman Bracket rules:
\begin{enumerate}
    \item The bracket of the unknot (no crossings) is $1$
        \[
            \langle \bigcirc \rangle = 1
        \]
    \item Replace a crossing with its two smoothings:
        \[
            \langle D \rangle = A \langle D_A \rangle + A^{-1} \langle D_B \rangle
        \]

        where $D_A$ is the diagram after A-smoothing, and $D_B$ is the diagram after B-smoothing.

    \item Removing a disjoint circle multiplies the result by an expression:
        \[
            \delta = -A^2-A^{-2}
        \]

        For example, the bracket of two separate circles is:

        \[
            -A^2-A^{-2}
        \]

        Adding another disjoint circle would give:

        \[
            (-A^2-A^{-2})^2
        \]

        Formally, adding another disjoint circle:

        \[
            \langle D \sqcup \bigcirc \rangle = \delta \langle D \rangle
        \]
\end{enumerate}

\subsubsection{Example}

Consider a knot diagram with a single crossing.

\hfill

Applying Rule 1 gives:

\[
    \langle D \rangle = A \langle D_A \rangle + A^{-1} \langle D_B \rangle
\]

Then, suppose the A-smoothing produces one circle, and the B-smoothing produces two circles.
We know:

\[
   \langle D_A \rangle = \langle \bigcirc \rangle = 1
\]

From rule 3, we also know:

\[
    \langle D_B \rangle = -A^2-A^{-2}
\]

Finally, from rule 2:

\[
   \langle D \rangle = A(1) + A^{-1}(-A^2-A^{-2})
\]

Expanding gives:

\[
    = A - A - A^{-3}
\]

Therefore:

\[
    \langle D \rangle = -A^{-3}
\]

\section{Implementing the Kauffman Bracket in a computer}

\subsection{Recursive implementation}

Since Rule 2 is already a recursive definition, it makes sense to implement the Kauffman Bracket in a computer recursively too.

\subsubsection{Basic algorithm}

A basic algorithm for calculating the Kauffman Bracket is:

\begin{enumerate}
    \item Receive a knot diagram as input via some data structure
    \item Count the number of crossings
    \item If there are no crossings, count the number of loops and return $\delta^{k-1}$, where $k$ is the number of disjoint circles.
    \item If there are still crossings, choose one crossing.
    \item Calculate the A-smoothing recursively
    \item Calculate the B-smoothing recursively
    \item Multiply the two results by $A$ and $A^{-1}$ respectively, then add them together.
\end{enumerate}

However, the issue with this is we still have not decided how to represent a knot diagram in a computer or how to calculate the two types of smoothings.

\end{document}
